\section{引言}
\subsection{问题背景}

大语言模型的训练需要大量计算，而一项训练任务能够使用的算力总是有限的。在预算 $C$ 内，扩大参数量 $N$ 与增加数据量 $D$ 都会增加开销，语料质量 $Q$ 和领域配比 $\bm p$ 又会改变训练效果。因此，合理配置资源需要回答：现有预算更适合用在哪一项上？标度律研究已给出规模与损失之间的经验关系，可据此比较参数和数据的投入\cite{kaplan2020scaling,hoffmann2022chinchilla,bahri2024explaining}。GPT-3 的少样本实验则从多类任务的表现说明了扩大规模的作用\cite{brown2020gpt3}。

继续扩大训练还面临语料供给问题。随着数据需求增长，高质量文本是否足够，成为“数据墙”研究关注的内容\cite{villalobos2022data}。与此同时，数据怎样组织、怎样使用，也值得考虑。2026 年《Nature》报道的 OpenScholar-8B 采用 80 亿参数，在引文准确性和减少幻觉等特定任务上优于 GPT-4o\cite{asai2026openscholar}。这一实例说明，分析资源配置时需要把质量和领域构成纳入讨论。语料处理、模型扩容和上下文延长各有费用，最终还要放到同一预算下比较。

围绕这些问题，题目提供了质量指标、配方实验、训练日志及能力评测等资料。本文利用这些记录，先研究参数、数据和质量怎样影响交叉熵损失，再计算预算内的配置；对于实际能力，则从历史评测出发分析增长来源，并估计未来可能达到的前沿。

\subsection{问题重述}

\textbf{问题1：} 整理 22 个质量指标，统一方向后给出文本、语料及领域三个层次的评分 $Q$；说明指标发生分歧时怎样识别和处理，并研究 17 个训练领域的占比与交叉熵损失之间的联系。

\textbf{问题2：} 将质量 $Q$、领域配比 $\bm{p}$ 纳入经典标度律，描述损失随 $N$、$D$、$Q$、$\bm{p}$ 的变化。利用边际效用和弹性比较不同投入，求出改善质量与扩大参数取得相同损失降幅的条件。

\textbf{问题3：} 把算力上限作为约束，安排参数量 $N$、数据量 $D$、配比 $\bm{p}$ 及质量 $Q$；当预算跨越若干数量级时，进一步判断主要投入方向会怎样变化。

\textbf{问题4：} 从历史能力记录中核算规模扩张与非规模技术进步所对应的增量及占比，并对开源大语言模型未来 12 至 24 个月的能力前沿作出估计。
